\section{Hoja de Observación del Flujo en el Acceso}
\apendiceindice{C}{Hoja de Observación del Flujo en el Acceso}

\subsection{Procedimiento}
La observación se realiza en el acceso piloto durante cinco días hábiles
consecutivos, desde 30 minutos antes hasta 10 minutos después del inicio de la
primera clase del turno vespertino, en intervalos de cinco minutos. Participan
dos observadores: el primero cuenta los ingresos y anota la longitud máxima de
la fila en cada intervalo; el segundo cronometra, para cinco estudiantes
elegidos al azar por intervalo, el tiempo desde que llegan al punto de revisión
hasta que lo cruzan. La misma hoja se utilizará durante el piloto para comparar
los resultados.

\subsection{Formato}
\noindent Fecha: \linea[3cm] \quad Día: \linea[2.5cm] \quad Acceso:
\linea[3cm]\\[0.5em]
Hora de inicio de clases: \linea[2cm] \quad Puntos de revisión activos ($c$):
\linea[1.5cm] \quad Observadores: \linea[3cm]

\medskip
\noindent
{\small
	\begin{tabular}{|p{2.2cm}|p{1.3cm}|p{1.3cm}|p{4.2cm}|p{1.2cm}|p{1.2cm}|p{1.2cm}|}
		\hline
		\textbf{Intervalo (min)}         & \textbf{Ingresos}  & \textbf{Fila máx.} &
		\textbf{Tiempos de revisión (s)} & \textbf{Sin cred.} & \textbf{Sin
		revisión}                        & \textbf{Rechaz.}                                    \\
		\hline
		$-30$ a $-25$                    &                    &                    &   &  &  & \\[0.8em] \hline
		$-25$ a $-20$                    &                    &                    &   &  &  & \\[0.8em] \hline
		$-20$ a $-15$                    &                    &                    &   &  &  & \\[0.8em] \hline
		$-15$ a $-10$                    &                    &                    &   &  &  & \\[0.8em] \hline
		$-10$ a $-5$                     &                    &                    &   &  &  & \\[0.8em] \hline
		$-5$ a $0$                       &                    &                    &   &  &  & \\[0.8em] \hline
		$0$ a $+5$                       &                    &                    &   &  &  & \\[0.8em] \hline
		$+5$ a $+10$                     &                    &                    &   &  &  & \\[0.8em] \hline
	\end{tabular}}

\subsection{Análisis}
Con los datos de cada intervalo se calculan la tasa de llegadas
$\lambda = \mathrm{ingresos}/5$ (estudiantes por minuto), el tiempo medio de
revisión $\bar{t}$ en segundos, la tasa de servicio de un punto de revisión
$\mu = 60/\bar{t}$ y la utilización del acceso:
\[
	\rho = \frac{\lambda}{c\,\mu}
\]
Cuando $\rho \geq 1$, llegan más estudiantes de los que el acceso puede atender
y la fila crece durante ese intervalo. Comparar $\bar{t}$, $\rho$ y la fila
máxima antes del piloto y durante él permite comprobar si el prototipo reduce
el tiempo de ingreso.
